% Consolidated 36-atomic-spectra from 3014fceb998eba4990e2b4a247494872c4c8e1cc; original TeX fragments preserved.

% === BEGIN 36-atomic-spectra-097-104.tex blob 7fefbb98e513fa97e87425598c4539aaed4fd60f ===
% Irodov Part VI §6.3, original 97--104, printed page 245.
% Input: user-provided PDF page 254 and appendix Rydberg constant.
\atomictermsproblem{97}{245}
锂原子的价电子在$2S$及$2P$态的结合能
分别为$5.39\,\mathrm{eV}$和$3.54\,\mathrm{eV}$。
求对应里德伯修正量。
\exsolution
碱金属结合能近似为
$E_{nL}=\mathrm{Ry}/(n+\alpha_L)^2$，
其中$\mathrm{Ry}=13.606\,\mathrm{eV}$。
本题$n=2$，故
\[
 \boxed{\alpha_S=\sqrt{13.606/5.39}-2\simeq-0.41},
 \qquad
 \boxed{\alpha_P=\sqrt{13.606/3.54}-2\simeq-0.04}.
\]
\excheck 采用$n+\alpha$记号故修正量为负，
常见量子亏损$\delta=-\alpha$为正。

\atomictermsproblem{98}{245}
钠价电子基态$3S$结合能为$5.14\,\mathrm{eV}$，
第一次激发电势为$2.10\,\mathrm V$。
求$3P$项的里德伯修正量。
\exsolution
激发态价电子结合能
$E_b(3P)=5.14-2.10=3.04\,\mathrm{eV}$，
故
\[
 \boxed{\alpha_{3P}=\sqrt{\mathrm{Ry}/3.04}-3
 \simeq-0.88}.
\]
\excheck 由于电子被激发后更容易电离，
$3P$结合能低于$3S$结合能。

\atomictermsproblem{99}{245}
锂锐线系第一条谱线波长$\lambda_1=813\,\mathrm{nm}$，
系列短波限$\lambda_2=350\,\mathrm{nm}$。
求基态价电子结合能。
\exsolution
锐线系对应$nS\to2P$。
短波极限给$E_{b,2P}=hc/\lambda_2$，
首条谱线$3S\to2P$给
$E_{b,3S}=hc(1/\lambda_2-1/\lambda_1)$。
S系列修正量
$\alpha_S=\sqrt{\mathrm{Ry}/E_{b,3S}}-3$，
所以
\[
 \boxed{E_{b,2S}=\frac{\mathrm{Ry}}{(2+\alpha_S)^2}
 \simeq5.3\,\mathrm{eV}}.
\]
\excheck 本题最终求的是$2S$基态结合能，
而非由系列短波限直接得出的$2P$结合能。

\atomictermsproblem{100}{245}
锂的S、P项里德伯修正量分别为$-0.41$、$-0.04$。
求由$3S$经$2P$至基态$2S$时发出的谱线波长。
\exsolution
电偶极跃迁要求$\Delta L=\pm1$，
故顺次发生$3S\to2P$、$2P\to2S$。
令$E_{b,nL}=\mathrm{Ry}/(n+\alpha_L)^2$，
\[
 \lambda_1=\frac{hc}{E_{b,2P}-E_{b,3S}},
 \qquad
 \lambda_2=\frac{hc}{E_{b,2S}-E_{b,2P}}.
\]
代入各$n,L$项，得
\[
 \boxed{\lambda_{3S\to2P}\simeq0.82\,\mu\mathrm m},
 \qquad
 \boxed{\lambda_{2P\to2S}\simeq0.68\,\mu\mathrm m}.
\]
\excheck 两光子能量的和等于初末能级差；
不能直接认为$3S\to2S$是一条允许电偶极谱线。

\atomictermsproblem{101}{245}
钠共振黄线双线的波长为
$589.00$及$589.56\,\mathrm{nm}$，
均对应$3P\to3S$，求$3P$项精细劈裂能。
\exsolution
因两线共用$3S$下能级，其光子能量之差
等于上能级的差：
\[
 \boxed{\Delta E=hc\left|
 \frac1{589.00\,\mathrm{nm}}-\frac1{589.56\,\mathrm{nm}}
 \right|\simeq2.0\times10^{-3}\,\mathrm{eV}}.
\]
\excheck 必须用倒波长差，而非波长的直接差。

\atomictermsproblem{102}{245}
铯锐线系第一线是双线，
波长为$1358.8$和$1469.5\,\mathrm{nm}$。
求后续谱线双线分量的角频率间隔。
\exsolution
锐线系同一末态$P$项存在恒定精细结构分裂，
因此全部锐线系双线角频率间隔一致。
\[
 \boxed{\Delta\omega=
 2\pi c\left|\frac1{1358.8\,\mathrm{nm}}
 -\frac1{1469.5\,\mathrm{nm}}\right|
 \simeq1.05\times10^{14}\,\mathrm{rad/s}}.
\]
\excheck 角频率比普通频率相差$2\pi$。

\atomictermsproblem{103}{245}
列出氢原子电子主量子数$n=3$时
所有可能的光谱项。
\exsolution
轨道量子数$l=0,1,2$，
电子自旋$s=1/2$，
耦合允许$j=l\pm1/2$，舍去负$j$。
于是
\[
 \boxed{3S_{1/2},\quad3P_{1/2},\quad3P_{3/2},\quad
 3D_{3/2},\quad3D_{5/2}}.
\]
\excheck $l=0$时只允许$j=1/2$，共五个谱项。

\atomictermsproblem{104}{245}
原子轨道角动量与自旋量子数$(L,S)$分别为
(a)$(3,2)$；(b)$(3,3)$；(c)$(2,5/2)$，
求总角动量量子数$J$的所有可能取值。
\exsolution
根据三角耦合规则，
$J=|L-S|,|L-S|+1,\ldots,L+S$，
\[
 \boxed{\text{(a)}\ J=1,2,3,4,5},
 \qquad
 \boxed{\text{(b)}\ J=0,1,2,3,4,5,6},
\]
\[
 \boxed{\text{(c)}\ J=\tfrac12,\tfrac32,\tfrac52,\tfrac72,\tfrac92}.
\]
\excheck 整数与半整数耦合时总$J$为半整数。
% === END 36-atomic-spectra-097-104.tex ===

% === BEGIN 36-atomic-spectra-105-117.tex blob 6de70a7cf1cdf23a809b451e7b703990a422d663 ===
% Irodov VI §6.3, 105--117. Printed pp.245--246; PDF pp.254--255.
% For term labels begin with \({}^{2S+1}L_J\), preserving JS-string safety.
\atomictermsproblem{105}{245}
求多重态\({}^{4}P\)及\({}^{5}D\)的所有可能总角动量模。
\exsolution
\({}^{4}P\)有$S=3/2,L=1$，故$J=1/2,3/2,5/2$；
\({}^{5}D\)有$S=2,L=2$，故$J=0,1,2,3,4$。
由$M_J=\hbar\sqrt{J(J+1)}$，
\[
 \boxed{{}^{4}P:\ \sqrt3\hbar/2,\ \sqrt{15}\hbar/2,\ \sqrt{35}\hbar/2},
\]
\[
 \boxed{{}^{5}D:\ 0,\sqrt2\hbar,\sqrt6\hbar,2\sqrt3\hbar,2\sqrt5\hbar}.
\]
\excheck 总角动量模是$\hbar\sqrt{J(J+1)}$，不是$J\hbar$。

\atomictermsproblem{106}{245}
求最大总角动量及相应谱项：
(a) Na价电子主量子数$n=4$；
(b) $1s^2$闭壳层外有$2p$与$3d$各一个电子。
\exsolution
(a)允许最大轨道$l=3$，单电子自旋$s=1/2$，
$J_{\max}=7/2$，
\[
 \boxed{{}^{2}F_{7/2},\quad M_{\max}=\sqrt{63}\hbar/2}.
\]
(b)两电子最大合轨道$L=1+2=3$，
平行自旋$S=1$，最大总$J=4$，
\[
 \boxed{{}^{3}F_4,\quad M_{\max}=\sqrt{20}\hbar}.
\]
\excheck 两未配对电子自旋平行产生三重态，
不能写成不可能的\({}^{2}F_4\)。

\atomictermsproblem{107}{245}
已知原子在F和D态的总量子数$J$
各有五种允许值，求两态的自旋角动量模。
\exsolution
耦合$J$不同取值数为$2\min(L,S)+1$。
F态$L=3$故$S=2$，给
\[
 \boxed{M_S(F)=\hbar\sqrt6}.
\]
D态$L=2$故仅可知$S\ge2$，给
\[
 \boxed{M_S(D)\ge\hbar\sqrt6}.
\]
\excheck D态并不能由五种$J$值唯一确定$S$。

\atomictermsproblem{108}{245}
某原子项多重度为3，
总角动量模为$\hbar\sqrt{20}$，
求可能的轨道量子数$L$。
\exsolution
多重度$2S+1=3$故$S=1$，
而$J(J+1)=20$故$J=4$。
角动量三角规则给
$|L-1|\le4\le L+1$，
\[
 \boxed{L=3,4,5}.
\]
\excheck 自旋量子数 $S=1$，总角动量 $J=4$；分别将 $L=3,4,5$ 代入三角不等式 $|L-S|\le J\le L+S$，三者均满足。它们依次对应 F、G、H 三个轨道项；其他整数轨道量子数均不满足。

\atomictermsproblem{109}{246}
求光谱项\({}^{x}D_2\)、
\({}^{x}P_{3/2}\)、\({}^{x}F_1\)
可能的多重度$x$。
\exsolution
由$x=2S+1$、$|L-S|\le J\le L+S$。
(a) $L=2,J=2$给$S=0,1,2,3,4$，
故$\boxed{x=1,3,5,7,9}$。
(b) $L=1,J=3/2$给$S=1/2,3/2,5/2$，
故$\boxed{x=2,4,6}$。
(c) $L=3,J=1$给$S=2,3,4$，
故$\boxed{x=5,7,9}$。
\excheck $J$为半整数时$S$也必须为半整数。

\atomictermsproblem{110}{246}
一个原子除闭壳层外有三电子，分别处于s、p、d轨道。
原子具有该构型最大总角动量。
求矢量模型中总自旋和总角动量夹角。
\exsolution
三自旋平行给$S=3/2$，
最大合轨道$L=0+1+2=3$，
最大$J=L+S=9/2$。
由$\boldsymbol J=\boldsymbol L+\boldsymbol S$，
\[
 \cos\theta=\frac{J(J+1)+S(S+1)-L(L+1)}
 {2\sqrt{J(J+1)S(S+1)}}.
\]
代入得到
\[
 \boxed{\theta=\arccos\left[
 \frac{99/4+15/4-12}{2\sqrt{(99/4)(15/4)}}\right]
 \simeq31^\circ}.
\]
\excheck 经典平行耦合图不能忽略量子角动量模的
$\sqrt{j(j+1)}$修正，故夹角并非0。

\atomictermsproblem{111}{246}
原子总角动量为$\hbar\sqrt6$，自旋$S=1$；
自旋与总角动量夹角$\theta=73.2^\circ$。
求光谱项。
\exsolution
由$J(J+1)=6$得$J=2$。
利用$\boldsymbol L=\boldsymbol J-\boldsymbol S$，
\[
 L(L+1)=J(J+1)+S(S+1)
 -2\cos\theta\sqrt{J(J+1)S(S+1)}\simeq6.
\]
故$L=2$为D项，多重度$2S+1=3$，
\[
 \boxed{{}^{3}D_2}.
\]
\excheck 要用三矢量夹角关系，不能直接用角度估计整数$L$。

\atomictermsproblem{112}{246}
由一个p电子和一个d电子组成的双电子系统，
求其所有允许的LS谱项。
\exsolution
两轨道耦合$L=1,2,3$；
自旋$1/2$与$1/2$耦合$S=0,1$。
单重态$J=L$；三重态$J=L-1,L,L+1$，
所以
\[
 \boxed{{}^{1}P_1,\ {}^{1}D_2,\ {}^{1}F_3,\
 {}^{3}P_{0,1,2},\ {}^{3}D_{1,2,3},\ {}^{3}F_{2,3,4}}.
\]
\excheck p和d为不等价电子，无等价壳层的额外泡利删项。

\atomictermsproblem{113}{246}
原子处于\({}^{2}P_{3/2}\)态，
再增加一个d电子，求合成体系可能的光谱项。
\exsolution
原有p态可取$L_1=1,S_1=1/2$，
新增d电子$L_2=2,S_2=1/2$，
重新LS耦合得$L=1,2,3,S=0,1$，
\[
 \boxed{{}^{1}P_1,\ {}^{1}D_2,\ {}^{1}F_3,\
 {}^{3}P_{0,1,2},\ {}^{3}D_{1,2,3},\ {}^{3}F_{2,3,4}}.
\]
\excheck 本题与前题的LS耦合项集合相同，
只是耦合次序的物理出发点不同。

\atomictermsproblem{114}{246}
判定以下电偶极跃迁的禁戒性：
(a) \({}^{2}D_{3/2}\to{}^{2}P_{1/2}\)；
(b) \({}^{3}P_1\to{}^{2}S_{1/2}\)；
(c) \({}^{3}F_3\to{}^{3}P_2\)；
(d) \({}^{4}F_{1/2}\to{}^{4}D_{5/2}\)。
\exsolution
选择规则$\Delta S=0$，$\Delta L=0,\pm1$，
$\Delta J=0,\pm1$（$0\leftrightarrow0$除外），
且宇称须改变。
(a)角动量条件允许。(b)多重度变3到2即$\Delta S\ne0$，
禁戒。(c)由F到P有$\Delta L=-2$，禁戒。
(d)按本版原题扫描，\({}^{4}F_{1/2}\)不满足
$L=3,S=3/2$下的$J\ge3/2$，
而所印$\Delta J=2$也禁戒。
故
\[
 \boxed{\text{(a)角动量允许；(b)(c)禁戒；
 (d)所印项本身不成立。}}
\]
\excheck 书末仅指出(b)(c)禁戒，
与(d)的印刷下标不一致；已保留争议而非静默改题。

\atomictermsproblem{115}{246}
求锂原子$3D$态的总简并度及其物理意义。
\exsolution
“$3D$”在此表示$n=3$、$L=2$，非三重态。
锂有一个价电子，$S=1/2$，
故$J=3/2,5/2$，
两子态简并度分别为4和6。
\[
 \boxed{g=4+6=10}.
\]
\excheck 未考虑精细结构时
也等于$(2L+1)(2S+1)=5\times2$。

\atomictermsproblem{116}{246}
求谱项\({}^{2}P\)、\({}^{3}D\)、\({}^{4}F\)
在其最大可能$J$下的简并度。
\exsolution
最大$J=L+S$，每个$J$下有$2J+1$个磁分量。
\[
 \boxed{g({}^{2}P)=2(1+\tfrac12)+1=4},\quad
 \boxed{g({}^{3}D)=2(2+1)+1=7},\quad
 \boxed{g({}^{4}F)=2(3+\tfrac32)+1=10}.
\]
\excheck 与整组未分裂项总简并度不同，
这里仅统计指定的最大$J$子项。

\atomictermsproblem{117}{246}
某光谱项简并度为7，且$L=3S$。
求该项符号。
\exsolution
$g=2J+1=7$给$J=3$。
由$L=3S$及$|L-S|\le J\le L+S$，
得$2S\le3\le4S$，
即$3/4\le S\le3/2$。
而$L=3S$必须为整数、$S$为整或半整数，
唯一可行$S=1,L=3$。
所以
\[
 \boxed{{}^{3}F_3}.
\]
\excheck 应联立三角耦合条件，不能只凭简并度判定$L$。
% === END 36-atomic-spectra-105-117.tex ===

% === BEGIN 36-atomic-spectra-118-134.tex blob 9d373b92eed99aaff8d78070ec530fba40b767a5 ===
% Irodov VI §6.3 original 118--134.
% Printed pages 246--247, supplied PDF pp.255--256, answers pp.345--346.
\atomictermsproblem{118}{246}
某元素原子的K、L、M三主壳层及$4s$副壳层均全满，
而$4p$副壳层恰半满，求元素名称与原子序数。
\exsolution
封闭K、L、M壳层电子数分别为$2,8,18$；
$4s$全满再贡献2个电子；
$4p$可容6个，半满即3个。
总电子数为
\[
 Z=2+8+18+2+3=\boxed{33}.
\]
故该元素为$\boxed{\mathrm{As}\text{（砷）}}$，
构型简记为$[\mathrm{Ar}]3d^{10}4s^24p^3$。
\excheck 题中M壳全满含$3d^{10}$，
不可只数$3s^23p^6$。

\atomictermsproblem{119}{246}
用洪特规则求只有一个未满$p$副壳层、
分别含(a)3个和(b)4个电子的原子基态谱项。
\exsolution
$p$有$l=1$及三个$m_l=-1,0,1$。
(a) $p^3$为半满，三个电子自旋平行，
$S=3/2$且$L=0$，故$J=3/2$，
\[
 \boxed{p^3:\ {}^{4}S_{3/2}}.
\]
(b) $p^4$等价于两个$p$空穴，
高自旋$S=1$，合轨道$L=1$；
超过半满时$J=L+S=2$，
\[
 \boxed{p^4:\ {}^{3}P_2}.
\]
\excheck 洪特第三规则：少于半满取$J=|L-S|$，
超过半满取$J=L+S$。

\atomictermsproblem{120}{246}
用洪特规则求仅有未满$d$副壳层且其电子数
分别为(a)3、(b)7时原子的基态总角动量模。
\exsolution
$d$轨道$l=2$，
对$d^3$有三自旋同向$S=3/2$，
最大$L=2+1+0=3$（F）。
少于半满取$J=|3-3/2|=3/2$，
\[
 \boxed{d^3:\ {}^{4}F_{3/2},\quad
 M_J=\frac{\sqrt{15}}2\hbar}.
\]
对$d^7$可用三空穴等效，
$S=3/2,L=3$；超过半满取
$J=3+3/2=9/2$，
\[
 \boxed{d^7:\ {}^{4}F_{9/2},\quad
 M_J=\frac{\sqrt{99}}2\hbar}.
\]
\excheck 总角动量模是$\hbar\sqrt{J(J+1)}$，
不是$J\hbar$；$d^7$按多于半满处理。

\atomictermsproblem{121}{246}
一个原子只有一个未满副壳层，基态分别为
(a)\({}^{3}F_2\)；
(b)\({}^{2}P_{3/2}\)；
(c)\({}^{6}S_{5/2}\)。
用洪特规则求该副壳层电子数。
\exsolution
(a) $F$态$L=3$、三重自旋$S=1$，
$J=|L-S|=2$表明少于半满；
可由两个$d$电子构成，
故$\boxed{d^2}$。
(b) $P$态$L=1,S=1/2$，
$J=L+S=3/2$，相应$p$超过半满，
故$\boxed{p^5}$。
(c) $S$态$L=0,S=5/2$，
来自五个自旋同向$d$电子的半满壳层，
故$\boxed{d^5}$。
\excheck 上述指唯一未满副壳层，
不包含闭合壳层内电子数。

\atomictermsproblem{122}{246}
用洪特规则确定唯一未满副壳层原子的基态：
(a)壳层装满了最大容量的$1/3$，且$S=1$；
(b)装满$70\%$，且$S=3/2$。
\exsolution
(a) 对$p$副壳层最大6个位置，
装入2个即占$1/3$，可得
$S=1,L=1$；少于半满取$J=0$，
\[
 \boxed{{}^{3}P_0}.
\]
(b) 对$d$副壳层最大10个位置，
装入7个即占$70\%$，其
$S=3/2,L=3$，超过半满取$J=9/2$，
\[
 \boxed{{}^{4}F_{9/2}}.
\]
\excheck 含电子数必须为整数，
同时满足指定的自旋量子数和填充分数。

\atomictermsproblem{123}{246}
一个原子仅有一个未满副壳层，
该副壳层有三电子，
其基态的轨道量子数$L=3$。
用洪特规则写出基态谱项。
\exsolution
$p^3$的最大轨道量子数为$L=0$，
因此$L=3$表明三电子应在$d$副壳层。
三个$d$电子自旋同向$S=3/2$，
最大轨道$L=3$。
少于半满，$J=|L-S|=3/2$，
\[
 \boxed{{}^{4}F_{3/2}}.
\]
\excheck 若先判定该电子构型为$d^3$，
再按洪特第三规则即可确定下标$J$。

\atomictermsproblem{124}{247}
一个原子的唯一未满副壳层恰含五电子且半满。
求其基态磁矩的模（以玻尔磁子表示）。
\exsolution
半满五电子只能为$d^5$，
洪特规则给$S=5/2,L=0$，
因而$J=5/2$，基态谱项为\({}^{6}S_{5/2}\)。
纯S态朗德因子$g=2$，
\[
 \boxed{\mu=g\mu_B\sqrt{J(J+1)}
 =2\mu_B\sqrt{(5/2)(7/2)}
 =\sqrt{35}\,\mu_B\simeq5.92\mu_B}.
\]
\excheck 此磁矩是向量模，不是沿外场的最大投影
$gJ\mu_B=5\mu_B$。

\atomictermsproblem{125}{247}
温度$T=3000\,\mathrm K$时，
有多少比例的氢原子处于主量子数$n=2$态？
假定稀薄气体满足热平衡且几乎全部在基态。
\exsolution
氢能级$E_n=-13.6/n^2\,\mathrm{eV}$，
故$1\to2$激发能为$\Delta E=10.2\,\mathrm{eV}$。
忽略微小精细结构，
$n$级统计权重与$n^2$成正比，
$g_2/g_1=4$。
玻尔兹曼比例
\[
 \boxed{\frac{N_2}{N_1}=
 4e^{-10.2\,\mathrm{eV}/(k_BT)}
 \simeq2.9\times10^{-17}}.
\]
\excheck 本题假设中性氢并处于热平衡；
如此极小的激发份额可近似为总原子中的比例。

\atomictermsproblem{126}{247}
钠蒸气温度$T=2400\,\mathrm K$。
求处于$3P$与基态$3S$的原子数比。
已知$3P\to3S$共振线波长$\lambda=589\,\mathrm{nm}$。
\exsolution
$3P$的轨道角动量$L=1,S=1/2$，
包括6个磁-自旋态；$3S$仅2个，
故$g_P/g_S=6/2=3$。
激发能$\Delta E=hc/\lambda$，
玻尔兹曼分布给
\[
 \boxed{\frac{N_{3P}}{N_{3S}}=
 3\exp\!\left[-\frac{hc}{\lambda k_BT}\right]
 \simeq1.14\times10^{-4}}.
\]
\excheck 必须保留简并度因子3，
不能单用$e^{-\Delta E/kT}$。

\atomictermsproblem{127}{247}
激发原子形成的定向束流速率$v=600\,\mathrm{m/s}$。
沿束流前进$l=2.5\,\mathrm{mm}$后，
通往基态的谱线强度降为原来的$1/\eta$，
其中$\eta=25$。
求平均激发态寿命。
\exsolution
激发态数目满足
$N(t)=N(0)e^{-t/\tau}$；
自发辐射强度正比于$N$。
飞行时间$t=l/v$，因而
$e^{-l/(v\tau)}=1/\eta$，
\[
 \boxed{\tau=\frac l{v\ln\eta}
 =\frac{2.5\times10^{-3}}
 {600\ln25}\,\mathrm s\simeq1.3\,\mu\mathrm s}.
\]
\excheck 沿程衰减是激发态统计指数衰减，
不是电磁波在介质中的比尔衰减。

\atomictermsproblem{128}{247}
稀薄汞蒸气原子本来几乎全部处于基态，
经共振波长$\lambda=253.65\,\mathrm{nm}$的汞灯照射。
蒸气以该波长再辐射功率$P=35\,\mathrm{mW}$，
共振态平均寿命$\tau=0.15\,\mu\mathrm s$。
求共振激发态原子数。
\exsolution
每个激发原子单位时间平均发射
$1/\tau$个共振光子，
光子能量$\varepsilon=hc/\lambda$。
所以总再辐射功率
$P=(N/\tau)(hc/\lambda)$，得到
\[
 \boxed{N=\frac{P\tau\lambda}{hc}
 \simeq6.7\times10^{9}\approx7\times10^9}.
\]
\excheck 已假定自发辐射主要经题给共振线退激；
若还有非辐射通道，则需要额外分支比。

\atomictermsproblem{129}{247}
锂原子浓度$n=3.6\times10^{16}\,\mathrm{cm^{-3}}$，
温度$T=1500\,\mathrm K$；
其$2P\to2S$共振线$\lambda=671\,\mathrm{nm}$
单位气体体积辐射功率为
$\mathcal P=0.30\,\mathrm{W/cm^3}$。
求共振激发态平均寿命。
\exsolution
热平衡激发态浓度
$n_{2P}=n(g_P/g_S)e^{-hc/(\lambda k_BT)}$，
其中$g_P/g_S=3$。
每个激发原子自发退激功率
$\varepsilon/\tau=hc/(\lambda\tau)$，
故
\[
 \mathcal P=
 \frac{3n\,e^{-hc/(\lambda k_BT)}}{\tau}
 \frac{hc}{\lambda}.
\]
解得
\[
 \boxed{\tau=
 \frac{3n\,hc}{\lambda\mathcal P}
 \exp\!\left[-\frac{hc}{\lambda k_BT}\right]
 \simeq6.6\times10^{-8}\,\mathrm s
 \approx65\,\mathrm{ns}}.
\]
\excheck $n$与$\mathcal P$都以$\mathrm{cm^{-3}}$
和$\mathrm{W/cm^3}$给出，体积单位在比值中抵消。

\atomictermsproblem{130}{247}
氢原子与自身辐射处于热平衡。
求：(a)$T=3000\,\mathrm K$时
从$2P$发射的受激与自发跃迁概率比；
(b)二者相等时的温度。
\exsolution
爱因斯坦系数结合普朗克辐射谱给
\[
 \frac{P_{\rm ind}}{P_{\rm sp}}=
 \frac1{\exp(\hbar\omega/k_BT)-1}.
\]
对$2P\to1S$共振跃迁，
$\hbar\omega=\Delta E=10.2\,\mathrm{eV}$。
(a) 代$T=3000\,\mathrm K$得
\[
 \boxed{\frac{P_{\rm ind}}{P_{\rm sp}}
 \simeq7.3\times10^{-18}}.
\]
(b) 比值等于1要求
$e^{\Delta E/(k_BT)}=2$，
\[
 \boxed{T=\frac{\Delta E}{k_B\ln2}
 \simeq1.71\times10^5\,\mathrm K}.
\]
\excheck 用户底本书末(a)印约$10^{-34}$，
与同一本书所用的$T=3000\,$K及
$10.2\,$eV代入相差约$10^{16}$倍。
本译解保留正确代入值，登记书末数值争议。
(b)的$1.71\times10^5$K是\emph{理想二能级
热平衡辐射模型的形式解}，不能据此断言
实际氢气仍由中性氢原子组成：
如此高温下电离平衡、电离连续谱以及
碰撞过程都必须计入，二能级近似不再自洽。

\atomictermsproblem{131}{247}
一束角频率等于气体原子两能级共振频率$\omega$的光，
通过温度$T$的稀薄原子气体。
考虑受激发射证明吸收系数
$\chi=\chi_0(1-e^{-\hbar\omega/k_BT})$，
其中$\hbar\omega\gg k_BT$，
$\chi_0$为低温极限吸收系数。
\exsolution
设低、上能级占据数密度分别$n_1,n_2$，
简并度$g_1,g_2$；
爱因斯坦系数满足$g_1B_{12}=g_2B_{21}$。
净吸收光功率与
$n_1B_{12}-n_2B_{21}$成正比，
因此
\[
 \chi=\frac{\hbar\omega}{c}
 (n_1B_{12}-n_2B_{21})
 =\frac{\hbar\omega}{c}n_1B_{12}
 \left(1-\frac{n_2g_1}{n_1g_2}\right).
\]
热平衡玻尔兹曼分布
$n_2/n_1=(g_2/g_1)e^{-\hbar\omega/k_BT}$，
且$\hbar\omega\gg k_BT$保证$n_1\simeq n$为总浓度，
所以
\[
 \boxed{\chi=\chi_0(1-e^{-\hbar\omega/k_BT})},
 \quad\chi_0=\frac{\hbar\omega nB_{12}}c.
\]
\excheck 该修正因子来自受激发射抵消一部分吸收，
并非改动单个原子的吸收系数定义。

\atomictermsproblem{132}{247}
汞共振线波长$\lambda=253.65\,\mathrm{nm}$，
激发态平均寿命$\tau=0.15\,\mu\mathrm s$。
求温度$T=300\,\mathrm K$时
多普勒展宽与自然线宽之比的数量级。
\exsolution
最概然热速率
$v_{\rm mp}=\sqrt{2k_BT/m_{\rm Hg}}$，
$m_{\rm Hg}\simeq200.59u$。
以$\Delta\omega_D\sim\omega_0v_{\rm mp}/c$
估计多普勒角频率宽度，
自然线宽$\Gamma\sim1/\tau$，
故
\[
 \frac{\Delta\omega_D}{\Gamma}
 \sim\frac{\omega_0v_{\rm mp}\tau}{c}
 =\frac{2\pi v_{\rm mp}\tau}{\lambda}.
\]
若对两侧速度展宽取全宽$2v_{\rm mp}$，
\[
 \boxed{\frac{\Delta\omega_{D,\mathrm{full}}}
 {\Delta\omega_{\mathrm{nat}}}
 \sim\frac{4\pi\tau v_{\rm mp}}{\lambda}
 \simeq1.2\times10^3\approx10^3}.
\]
\excheck 展宽的标准差、半高全宽与粗估全宽
有数量级为1的系数差；
书末亦仅要求估计$\sim10^3$。

\atomictermsproblem{133}{247}
铁（$Z=26$）的$K_\alpha$特征线波长为
$\lambda_{\rm Fe}=193\,\mathrm{pm}$。
利用莫塞莱定律求铜（$Z=29$）
的$K_\alpha$线波长。
\exsolution
对轻中等原子序数的$K_\alpha$线，
$\omega_\alpha\propto(Z-1)^2$，
故$\lambda_\alpha\propto(Z-1)^{-2}$。
于是
\[
 \boxed{\lambda_{\rm Cu}=
 \lambda_{\rm Fe}\left(\frac{26-1}{29-1}\right)^2
 =193\,\mathrm{pm}\frac{625}{784}
 \simeq154\,\mathrm{pm}}.
\]
\excheck 不能直接用$Z^{-2}$取代
题意所用含屏蔽修正的$(Z-1)^{-2}$。

\atomictermsproblem{134}{247}
利用莫塞莱定律求：
(a)铝和钴的$K_\alpha$线波长；
(b)钒原子K、L层电子结合能之差。
\exsolution
类氢屏蔽近似对K线取有效核电荷$Z-1$，
则
\[
 \Delta E_{K\alpha}=
 \mathrm{Ry}(Z-1)^2\left(1-\frac14\right)
 =\frac34\mathrm{Ry}(Z-1)^2,
 \quad \lambda_{K\alpha}=\frac{hc}{\Delta E_{K\alpha}}.
\]
对铝$Z=13$、钴$Z=27$，
\[
 \boxed{\lambda_{\rm Al}\simeq843\,\mathrm{pm}},\qquad
 \boxed{\lambda_{\rm Co}\simeq180\,\mathrm{pm}}.
\]
钒$Z=23$时，K、L结合能差即
$\Delta E_{K\alpha}$，因此
\[
 \boxed{E_K-E_L=
 \frac34(13.606\,\mathrm{eV})(22)^2
 \simeq4.94\,\mathrm{keV}\approx5\,\mathrm{keV}}.
\]
\excheck 该模型忽略了精细结构与精密屏蔽常数，
适用于本题要求的莫塞莱数量级。
% === END 36-atomic-spectra-118-134.tex ===

% === BEGIN 36-atomic-spectra-135-147.tex blob 98d10226671de3c51d3c8c6d5f589bc48a47752c ===
% Irodov VI §6.3 originals 135--147, printed p248, PDF p257.
% Answers user PDF pages 354--355.
\atomictermsproblem{135}{248}
两种元素的特征$K_\alpha$线波长分别为
$250\,\mathrm{pm}$和$179\,\mathrm{pm}$。
求它们在周期表中间隔着多少个元素。
\exsolution
莫塞莱定律给
$\lambda_{K\alpha}=hc/[(3/4)\mathrm{Ry}(Z-1)^2]$，
所以
\[
 Z=1+\sqrt{\frac{hc}{(3/4)\mathrm{Ry}\lambda}}.
\]
对$250\,\mathrm{pm}$得$Z_1\simeq23$（钒），
对$179\,\mathrm{pm}$得$Z_2\simeq27$（钴）。
中间的元素序数为24、25、26，故
\[
 \boxed{N=Z_2-Z_1-1=3}.
\]
\excheck “位于两者之间”不包含端点元素，
因而是三种而不是四种。

\atomictermsproblem{136}{248}
以镍（$Z=28$）为靶的X射线管，
其$K_\alpha$谱线与连续谱短波限的波长差
为$\Delta\lambda=84\,\mathrm{pm}$。
求加速电压$U$。
\exsolution
莫塞莱近似给
$K_\alpha$光子能量
$\Delta E_\alpha=(3/4)\mathrm{Ry}(27)^2$，
故其波长$\lambda_\alpha=hc/\Delta E_\alpha
\simeq167\,\mathrm{pm}$。
连续谱短波限更短，
$\lambda_{\min}=\lambda_\alpha-\Delta\lambda
\simeq83\,\mathrm{pm}$。
由$eU=hc/\lambda_{\min}$，
\[
 \boxed{U=\frac{hc}{e(\lambda_\alpha-\Delta\lambda)}
 \simeq15\,\mathrm{kV}}.
\]
\excheck 若误把短波限设为$\lambda_\alpha+\Delta\lambda$，
会得到不合理的低电压；短波限应位于更短波一侧。

\atomictermsproblem{137}{248}
铝靶X射线管的连续谱短波限
为$\lambda_{\min}=0.50\,\mathrm{nm}$，
而其K系激发电压为$U_K=1.56\,\mathrm{kV}$。
问这时还能观察到特征K系谱线吗？
\exsolution
连续谱短波限对应最高电子动能
$eU=hc/\lambda_{\min}$。
因此
\[
 U=\frac{hc}{e(0.50\,\mathrm{nm})}
 \simeq2.48\,\mathrm{kV}.
\]
因$2.48\,\mathrm{kV}>1.56\,\mathrm{kV}$，
入射电子足以从K层击出电子，
随后更高层电子填空位便发出K系特征线，
\[
 \boxed{\text{可以观察到K系谱线。}}
\]
\excheck 产生连续谱不自动保证能产生某一特征系，
要比较实际电子能量与对应内层电离阈值。

\atomictermsproblem{138}{248}
某X射线管的电压从$U_1=10\,\mathrm{kV}$
升高为$U_2=20\,\mathrm{kV}$。
底本将所比较的固定特征波长称为$K_\infty$线，
并说它与连续谱短波限的波长间距
增为原来的$\eta=3.0$倍。
求靶元素原子序数$Z$，并核查底本答案。
\exsolution
先单独解决\emph{几何上的波长间距}问题。
两电压下连续谱短波限
\[
 \lambda_i=\frac{hc}{eU_i},\qquad
 \lambda_1=123.984\,\mathrm{pm},\quad
 \lambda_2=61.992\,\mathrm{pm}.
\]
令待定特征波长为$\lambda_X$。
若初始$10$kV时该特征线已经能被激发，
必须有$\lambda_X\ge\lambda_1$，
所以波长间隔满足
\[
 \lambda_X-\lambda_2
 =\eta(\lambda_X-\lambda_1),\qquad
 \boxed{\lambda_X=
 \frac{\eta\lambda_1-\lambda_2}{\eta-1}
 =154.980\,\mathrm{pm}}.
\]
它对应的光子能量为
$\varepsilon_X=hc/\lambda_X=8.00\,\mathrm{keV}$。
\emph{关键是辨认特征线的种类：}

若按底本书末采用的\emph{莫塞莱$K_\alpha$线}
近似能量
$\varepsilon_{K_\alpha}
=\frac34\mathrm{Ry}(Z-1)^2$，则
\[
 \boxed{Z=1+\sqrt{\frac{4\varepsilon_X}{3\mathrm{Ry}}}
 \simeq29.0},\qquad
 \boxed{\text{靶材为铜（Cu）}}.
\]
这与书末$Z=29$及其所列含$3\mathrm{Ry}/4$
的公式相符。
但原题\emph{字面写$K_\infty$}，
它是K系极限（对应K层电离边），
最简氢样屏蔽模型应使用
$\varepsilon_{K_\infty}\simeq\mathrm{Ry}(Z-1)^2$，
因而
\[
 \boxed{Z_{K_\infty}\simeq
 1+\sqrt{\frac{8.00\times10^3}{13.606}}
 \simeq25.25\approx25}
\]
（模型精度内对应锰），而非29。
再者，若把波长“间距”理解为绝对值且
\emph{没有}要求第一电压下产生特征线，
还存在跨越两短波限的代数分支：
\[
 |\lambda_X-\lambda_2|=3|\lambda_X-\lambda_1|,
 \quad
 \lambda_2<\lambda_X<\lambda_1
 \ \Longrightarrow\
 \lambda_X=\frac{3\lambda_1+\lambda_2}{4}
 =108.486\,\mathrm{pm}.
\]
此时$\varepsilon_X=11.43$keV，
超过第一电压电子最高能量10keV，
第一状态不能激发该线。
\excheck 已对底本原题图像和书末第138题公式直接校核：
\emph{原题写$K_\infty$，书末解法却用了
$K_\alpha$的$\frac34$莫塞莱系数}。
所以保留\emph{书末口径}的$Z=29$，
同时列出严格按字面$K_\infty$解读的
$Z\approx25$与未激发分支；
不能仅用前版“书末答案错误”一句带过，
也不可把$K_\infty$静默改译成$K_\alpha$。

\atomictermsproblem{139}{248}
某金属X射线K、L两吸收限对应角频率之差为
$\Delta\omega=6.85\times10^{18}\,\mathrm{s^{-1}}$。
求该金属。
\exsolution
在莫塞莱模型中，
K、L结合能的差约为
$\Delta E=(3/4)\mathrm{Ry}(Z-1)^2$，
故
\[
 \Delta\omega=\frac{\Delta E}{\hbar}
 =\frac{3\mathrm{Ry}}{4\hbar}(Z-1)^2.
\]
反解
\[
 \boxed{Z=1+\sqrt{\frac{4\hbar\Delta\omega}
 {3\mathrm{Ry}}}\simeq22}.
\]
原子序数22是$\boxed{\mathrm{Ti}\text{（钛）}}$。
\excheck 题目给的是\emph{角频率}，
因此用$\Delta E=\hbar\Delta\omega$，不是$h\Delta\omega$。

\atomictermsproblem{140}{248}
钒（$Z=23$）的L层吸收限波长
$\lambda_L=2.4\,\mathrm{nm}$。
求其K层电子结合能。
\exsolution
L层结合能为
$E_L=hc/\lambda_L\simeq0.517\,\mathrm{keV}$。
莫塞莱近似给K、L层结合能差
$\Delta E_{KL}=(3/4)\mathrm{Ry}(Z-1)^2$，
对钒$Z-1=22$，
$\Delta E_{KL}\simeq4.94\,\mathrm{keV}$。
所以
\[
 \boxed{E_K=E_L+\Delta E_{KL}\simeq5.46\,\mathrm{keV}
 \approx5.5\,\mathrm{keV}}.
\]
\excheck 不能直接用K$\alpha$能量当K电子全部结合能，
因为该谱线对应两层能量\emph{之差}。

\atomictermsproblem{141}{248}
已知钛的K系第一条谱线$K_\alpha$
与本系短波极限的波长差为
$\Delta\lambda=26\,\mathrm{pm}$。
求钛L层电子结合能。
\exsolution
钛$Z=22$。
莫塞莱关系给K$\alpha$能量
$E_\alpha=(3/4)\mathrm{Ry}(21)^2
\simeq4.50\,\mathrm{keV}$，
所以$\lambda_\alpha=hc/E_\alpha\simeq275.5\,\mathrm{pm}$。
K系短波极限对应K层结合能，
其波长
$\lambda_K=\lambda_\alpha-\Delta\lambda
\simeq249.5\,\mathrm{pm}$。
于是
$E_K=hc/\lambda_K\simeq4.97\,\mathrm{keV}$，
$E_L=E_K-E_\alpha$，故
\[
 \boxed{E_L\simeq0.47\,\mathrm{keV}
 \approx0.5\,\mathrm{keV}}.
\]
\excheck 两谱线间隔是波长差，不等于能量差；
计算必须分别取$hc/\lambda$再相减。

\atomictermsproblem{142}{248}
锌的$K_\alpha$辐射照射铁，
从铁K层打出电子。
已知铁的K吸收限波长$\lambda_K=174\,\mathrm{pm}$。
求光电子动能和速率。
\exsolution
锌$Z=30$的K$\alpha$光子能量
$\varepsilon=(3/4)\mathrm{Ry}(29)^2
\simeq8.58\,\mathrm{keV}$。
铁K层结合能
$E_{b,\mathrm{Fe}}=hc/(174\,\mathrm{pm})
\simeq7.13\,\mathrm{keV}$。
光电子动能
\[
 \boxed{T=\varepsilon-E_b\simeq1.45\,\mathrm{keV}}.
\]
在此能区可用非相对论近似
$v=\sqrt{2T/m_e}$，故
\[
 \boxed{v\simeq2.26\times10^7\,\mathrm{m/s}}.
\]
\excheck 必须用锌的特征谱线能量减去\emph{铁}的K层结合能。

\atomictermsproblem{143}{248}
求朗德$g$因子：(a) S项；(b)单重态。
\exsolution
一般LS耦合有
\[
 g_J=1+\frac{J(J+1)+S(S+1)-L(L+1)}
 {2J(J+1)}.
\]
(a) S项$L=0,J=S>0$，
故$\boxed{g_J=2}$。
若同时$S=0$则$J=0$，上述公式为$0/0$，
磁矩本身为零。
(b) 单重态$S=0$、$J=L>0$，
故$\boxed{g_J=1}$。
\excheck 单重S态的$J=0$使$g$的常规定义不适用，
不能强行把$0/0$写成某个数。

\atomictermsproblem{144}{248}
求下列光谱项的朗德因子：
(a)\({}^{6}F_{1/2}\)；
(b)\({}^{4}D_{1/2}\)；
(c)\({}^{5}F_2\)；
(d)\({}^{5}P_1\)；
(e)\({}^{3}P_0\)。
\exsolution
代入
$g=1+[J(J+1)+S(S+1)-L(L+1)]/[2J(J+1)]$。
(a)$(S,L,J)=(5/2,3,1/2)$，故$g=-2/3$；
(b)$(3/2,2,1/2)$，得$g=0$；
(c)$(2,3,2)$，得$g=1$；
(d)$(2,1,1)$，得$g=5/2$；
(e)$J=0$，公式分母为零、磁矩为零，
$g$不作通常定义。于是
\[
 \boxed{g_a=-\tfrac23,\quad g_b=0,\quad
 g_c=1,\quad g_d=\tfrac52,\quad g_e:\text{未定义}}.
\]
\excheck $g$可为负，零总角动量态不应硬算$g$的数值。

\atomictermsproblem{145}{248}
求以下原子磁矩模，以玻尔磁子为单位：
(a)单重F态；
(b)\({}^{2}D_{3/2}\)；
(c)$S=1,L=2,g_J=4/3$的状态。
\exsolution
取
$\mu=g_J\mu_B\sqrt{J(J+1)}$。
(a) $S=0,L=J=3,g=1$，
\[
 \boxed{\mu_a=\sqrt{12}\,\mu_B=2\sqrt3\,\mu_B}.
\]
(b) $S=1/2,L=2,J=3/2$，
$g=1+[(15/4)+(3/4)-6]/[2(15/4)]=4/5$，故
\[
 \boxed{\mu_b=\frac45\sqrt{\frac{15}4}\,\mu_B
 =\frac{2\sqrt{15}}5\,\mu_B}.
\]
(c) $S=1,L=2$，由$g_J=4/3$反求$J=3$，
因此
\[
 \boxed{\mu_c=\frac43\sqrt{12}\,\mu_B
 =\frac{8\sqrt3}{3}\,\mu_B}.
\]
\excheck (c)必须先从$g$确定$J$，
不能任取$J=1$或2。

\atomictermsproblem{146}{248}
一个原子的谱项为$D_2$，
磁矩沿外场的最大投影为$4\mu_B$。
求其总自旋角动量模。
\exsolution
$D$给$L=2$，下标给$J=2$。
最大磁矩投影$g_JJ\mu_B=4\mu_B$，
故$g_J=2$。
朗德公式化简为
\[
 g_J=1+\frac{S(S+1)}{2J(J+1)}
 =1+\frac{S(S+1)}{12}.
\]
令$g_J=2$得$S(S+1)=12$，
即$S=3$，所以
\[
 \boxed{M_S=\hbar\sqrt{S(S+1)}=
 2\sqrt3\,\hbar}.
\]
\excheck 此题使用的是磁矩\emph{沿场最大投影}，
不是磁矩模$g\sqrt{J(J+1)}\mu_B$。

\atomictermsproblem{147}{248}
某原子$L=2,S=1$，处于弱磁场中。
已知总角动量与磁场方向的最小夹角为$30^\circ$。
求原子磁矩的模。
\exsolution
空间量子化下总角动量最大投影
$M_{J,z}=J\hbar$，而模为
$M_J=\hbar\sqrt{J(J+1)}$，
故最小夹角满足
\[
 \cos\theta_{\min}=
 \frac J{\sqrt{J(J+1)}}=
 \sqrt{\frac J{J+1}}.
\]
取$\theta_{\min}=30^\circ$，
解得$J=3$。
对$L=2,S=1,J=3$，
$g_J=4/3$，因此
\[
 \boxed{\mu=\frac43\mu_B\sqrt{3\cdot4}
 =\frac{8\sqrt3}{3}\,\mu_B}.
\]
\excheck 磁矩模与沿磁场投影不同；
前者有$\sqrt{J(J+1)}$系数。
% === END 36-atomic-spectra-135-147.tex ===

% === BEGIN 36-atomic-spectra-148-158.tex blob 6a9921545825bb223cebb80398055dddfbc22cb1 ===
% Irodov VI §6.3 original 148--158, printed p249, PDF p258.
% Note: source problem 155 prints "22 MeV", answer presumes 22 meV.
\atomictermsproblem{148}{249}
钠原子价电子处于主量子数$n=3$的状态，
且总角动量达到该主量子数所允许的最大值。
求原子磁矩模。
\exsolution
$n=3$允许最大$l=2$（D态），电子自旋$s=1/2$，
所以最大$J=l+s=5/2$，
谱项\({}^{2}D_{5/2}\)。
朗德因子
\[
 g_J=1+\frac{J(J+1)+s(s+1)-l(l+1)}
 {2J(J+1)}
 =\frac65.
\]
因此
\[
 \boxed{\mu=g_J\mu_B\sqrt{J(J+1)}
 =\frac65\mu_B\frac{\sqrt{35}}2
 =\frac{3\sqrt{35}}5\mu_B}.
\]
\excheck 题中$n=3$仍可有$3d$轨道，
最大$J$不是$3p$态的$3/2$。

\atomictermsproblem{149}{249}
某激发原子的电子构型为$1s^22s^22p\,3d$，
且处于该构型最大总角动量态。
求原子磁矩的大小。
\exsolution
闭合壳层角动量相互抵消；
未闭合$p,d$各一电子，
取最大$L=1+2=3$，
两个自旋平行$S=1$，最大$J=L+S=4$，
即\({}^{3}F_4\)。
朗德因子
\[
 g_J=1+\frac{20+2-12}{2(20)}=\frac54,
\]
所以
\[
 \boxed{\mu=\frac54\sqrt{20}\,\mu_B
 =\frac{5\sqrt5}2\,\mu_B}.
\]
\excheck $J$最大不代表$g=2$，
轨道与自旋磁旋比分别不同。

\atomictermsproblem{150}{249}
原子处于$S=3/2,L=2$态，
已知其总磁矩为零。
求总角动量模。
\exsolution
非零$J$时磁矩模
$\mu=g_J\mu_B\sqrt{J(J+1)}$。
允许$J=1/2,3/2,5/2,7/2$。
由$g_J=0$得
\[
 1+\frac{J(J+1)+15/4-6}{2J(J+1)}=0,
\quad3J(J+1)=\frac94.
\]
所以$J(J+1)=3/4$即$J=1/2$，
\[
 \boxed{M_J=\frac{\sqrt3}2\hbar}.
\]
\excheck 此处磁矩为零来自朗德因子$g_J=0$，
并非总角动量矢量模为零。

\atomictermsproblem{151}{249}
某原子自旋$S=2$，总角动量模为$\sqrt2\,\hbar$，
且磁矩为零。写出相应光谱项。
\exsolution
由$J(J+1)=2$得$J=1$。
磁矩零要求$g_J=0$，代入朗德公式：
\[
 1+\frac{2+S(S+1)-L(L+1)}{2(2)}=0.
\]
$S(S+1)=6$，整理得$L(L+1)=12$，
因此$L=3$（F项）。
多重度$2S+1=5$，
\[
 \boxed{{}^{5}F_1}.
\]
\excheck $S=2,L=3$允许$J=1$，
符合角动量三角不等式。

\atomictermsproblem{152}{249}
处于\({}^{2}P_{3/2}\)态的原子
置于磁感应强度$B=1.0\,\mathrm{kG}$的均匀磁场。
求总角动量矢量绕磁场的进动角频率。
\exsolution
对\({}^{2}P_{3/2}\)，$L=1,S=1/2,J=3/2$，
朗德因子
$g_J=1+[15/4+3/4-2]/[2(15/4)]=4/3$。
弱场进动频率
$\omega_L=g_J\mu_B B/\hbar$，
$1\,\mathrm{kG}=0.10\,\mathrm T$，
所以
\[
 \boxed{\omega_L=
 \frac{4\mu_B(0.10\,\mathrm T)}{3\hbar}
 \simeq1.17\times10^{10}\,\mathrm{rad/s}}.
\]
\excheck 高斯单位的$1$kG须换算为$0.1$T，
不可数值原样代入SI公式。

\atomictermsproblem{153}{249}
\({}^{2}P_{1/2}\)态原子在圆环电流轴上，
圆环半径$r=5.0\,\mathrm{cm}$，
电流$I=10\,\mathrm A$，原子离环心的轴向距离也为$r$。
求该环对原子的最大磁力。
\exsolution
\({}^{2}P_{1/2}\)有$L=1,S=1/2,J=1/2$，
$g_J=2/3$；
最大磁矩沿轴投影
$|\mu_z|=g_JJ\mu_B=\mu_B/3$。
SI制圆环轴向磁场
\[
 B(z)=\frac{\mu_0Ir^2}{2(r^2+z^2)^{3/2}},
 \qquad
 \left|\frac{dB}{dz}\right|_{z=r}
 =\frac{3\mu_0I}{2^{7/2}r^2}.
\]
因此
\[
 \boxed{F_{\max}=|\mu_z B'|
 =\frac{\mu_0I\mu_B}{2^{7/2}r^2}
 \simeq4.1\times10^{-27}\,\mathrm N}.
\]
\excheck 应使用磁矩的\emph{轴向投影}；
不能将磁矩模代入轴向梯度力。

\atomictermsproblem{154}{249}
基态氢原子在距长直导线$r=2.5\,\mathrm{cm}$处，
导线电流$I=10\,\mathrm A$。
求原子所受磁力的最大值。
\exsolution
氢基态为\({}^{2}S_{1/2}\)，$g_J=2,J=1/2$，
最大磁矩沿场方向投影
$|\mu_B|$。
长直导线场
$B(r)=\mu_0I/(2\pi r)$，
梯度模
$|dB/dr|=\mu_0I/(2\pi r^2)$。
因此
\[
 \boxed{F_{\max}=\mu_B\frac{\mu_0 I}{2\pi r^2}
 \simeq3.0\times10^{-26}\,\mathrm N}.
\]
\excheck 实际方向由原子的$m_J$取值决定，
本题给出的数值是力的最大模。

\atomictermsproblem{155}{249}
窄钒原子束处于基态\({}^{4}F_{3/2}\)，
经过长$l_1=5.0\,\mathrm{cm}$的横向非均匀磁场，
屏位于磁场后$l_2=15\,\mathrm{cm}$，
两个最外侧分束在屏上相距$\delta=2.0\,\mathrm{mm}$。
原印动能是$T=22\,\mathrm{MeV}$。
求磁场梯度，并说明与书末答案的差异。
\exsolution
该态$L=3,S=3/2,J=3/2$，
$g_J=1+[15/4+15/4-12]/[2(15/4)]=2/5$。
两最外侧$m_J=\pm J$磁矩投影差为
$\Delta\mu=2g_JJ\mu_B$。
设原子速度$v$，
磁区内时间$t_1=l_1/v$，
离开磁区至屏时间$t_2=l_2/v$，
则两束位移差
\[
 \delta=\frac{\Delta\mu}{M}B'
 \left(\frac{t_1^2}2+t_1t_2\right).
\]
用$T=Mv^2/2$消去质量和速率，
\[
 \boxed{B'=\frac{2T\delta}
 {g_JJ\mu_B l_1(l_1+2l_2)}}.
\]
若\emph{严格照原印}$T=22\,\mathrm{MeV}$，
\[
 \boxed{B'\simeq1.45\times10^{11}\,\mathrm{T/m}}.
\]
若能量单位实际应为$22\,\mathrm{meV}$，
则
\[
 \boxed{B'\simeq1.45\times10^2\,\mathrm{T/m}
 \simeq15\,\mathrm{kG/cm}},
\]
与原书书末给出的15kG/cm一致。
\excheck 原印MeV与书末数值隐含的meV相差
$10^9$倍；本译本明确列出两种结果，
不未经授权替换题面单位。

\atomictermsproblem{156}{249}
以下谱项在弱磁场中分裂为多少个子能级：
(a)\({}^{3}P_0\)；
(b)\({}^{2}F_{5/2}\)；
(c)\({}^{4}D_{1/2}\)？
\exsolution
线性塞曼能移
$\Delta E=g_J m_J\mu_BB$；
$m_J=-J,\ldots,J$。
(a)$J=0$，只有$m_J=0$，不劈裂；
(b)$J=5/2$且$g_J\ne0$，共有
$2J+1=6$个不同子能级；
(c)$L=2,S=3/2,J=1/2$，
朗德因子
\[
 g_J=1+\frac{3/4+15/4-6}{2(3/4)}=0,
\]
两磁态不发生一阶能量分离。
因此
\[
 \boxed{\text{(a)不分裂；(b)6个；(c)一阶不分裂}}.
\]
\excheck “(c)不分裂”不代表不存在两个$m_J$态，
而是它们在一阶弱场中能量重合。

\atomictermsproblem{157}{249}
弱磁场$B=2.50\,\mathrm{kG}$使谱项发生塞曼劈裂。
求(a)单重D项；(b)\({}^{3}F_4\)
的全劈裂能宽（eV）。
\exsolution
对总角动量$J$的子项，
端到端全劈裂
$\Delta E_{\rm tot}=2J|g_J|\mu_BB$。
(a)单重D项$L=J=2,S=0,g_J=1$，
\[
 \boxed{\Delta E_a=4\mu_B(0.25\,\mathrm T)
 \simeq5.79\times10^{-5}\,\mathrm{eV}}.
\]
(b)\({}^{3}F_4\)有$L=3,S=1,J=4$，
$g_J=5/4$，
\[
 \boxed{\Delta E_b=8(5/4)\mu_B(0.25\,\mathrm T)
 \simeq1.45\times10^{-4}\,\mathrm{eV}}.
\]
\excheck 这里全宽包含最大$m_J=+J$与最小$m_J=-J$
两端之差，不是相邻子级间隔。

\atomictermsproblem{158}{249}
在弱磁场中判断下列光谱跃迁呈正常还是反常塞曼效应：
(a)\({}^{1}P\to{}^{1}S\)；
(b)\({}^{2}D_{5/2}\to{}^{2}P_{3/2}\)；
(c)\({}^{3}D_1\to{}^{3}P_0\)；
(d)\({}^{5}I_5\to{}^{5}H_4\)。
\exsolution
线频移取决于
$\Delta\omega=(g_um_u-g_lm_l)\mu_BB/\hbar$，
仅当允许跃迁的频移组成等距三线群，
才表现为正常塞曼效应。
(a)单重态$S=0$、$g=1$，正常。
(b)两态分别$g_u=6/5$、
$g_l=4/3$，频移不是单一等间隔三线群，反常。
(c)下态$J_l=0$，上态$J_u=1$且
$g_u=1/2$，
只给三个等间隔频移，正常。
(d)\({}^{5}I_5\)有$L=6,S=2,J=5$，
\({}^{5}H_4\)有$L=5,S=2,J=4$，
分别计算朗德因子均为$9/10$，
所以允许$\Delta m=0,\pm1$只产生三线群，正常。
\[
 \boxed{\text{(a)正常；(b)反常；(c)正常；(d)正常}}.
\]
\excheck “正常”指观察到的等间隔三线型，
不要求所有项都恰有$g=1$。
% === END 36-atomic-spectra-148-158.tex ===

% === BEGIN 36-atomic-spectra-159-166.tex blob 1ef2f1cf7633a83980d525043c5c9fed83fb3edd ===
% Irodov VI §6.3 original 159--166 (end of §6.3).
% Printed p250, user-provided PDF p259, answers PDF p355.
\atomictermsproblem{159}{250}
原子某一\emph{单重谱项}在弱磁场
$B=3.0\,\mathrm{kG}$中全劈裂能宽，
本版原印为$\Delta E=10^{4}\,\mu\mathrm{eV}$。
求其光谱项符号，并核对原书书末所列结果。
\exsolution
单重态$S=0$，所以$J=L$且$g_J=1$，
全劈裂能宽为
\[
 \Delta E=2J\mu_BB=2L\mu_BB.
\]
由于$B=0.300\,\mathrm T$，
$\mu_B B\simeq17.365\,\mu\mathrm{eV}$。
\emph{逐字按题面}$10^4\,\mu\mathrm{eV}$
代入得
\[
 L=\frac{10^4}{2\times17.365}\simeq288,
\]
即
\[
 \boxed{L\simeq288\quad
 (\text{不对应题目通常涉及的低轨道原子项})}.
\]
原书书末答案写为\({}^{1}F_3\)；
这一答案要求$L=3$，
相应全宽却是
\[
 \boxed{\Delta E_{L=3}=6\mu_BB
 \simeq104.2\,\mu\mathrm{eV}}.
\]
\excheck 原印$10^{4}\,\mu\mathrm{eV}$
和书末\({}^{1}F_3\)之间差约96倍，
可能是排版将$104$错误设成$10^4$，
但本稿不擅自将原题改为$104\,\mu\mathrm{eV}$。
此题标记为\emph{原书印刷差异待核验}。

\atomictermsproblem{160}{250}
某一由单重项间跃迁形成的光谱线
波长$\lambda=612\,\mathrm{nm}$。
在$B=10.0\,\mathrm{kG}$的磁场中，
求最外侧两塞曼分量的波长间隔。
\exsolution
单重项$g=1$，正常塞曼效应中
$\Delta m=0,\pm1$对应频移
$0,\pm\mu_B B/\hbar$。
两端的角频率总差
$\Delta\omega_{\rm ext}=2\mu_BB/\hbar$。
小频移满足
$|\Delta\lambda|=(\lambda^2/2\pi c)|\Delta\omega|$，
所以
\[
 \boxed{\Delta\lambda_{\rm ext}=
 \frac{\lambda^2\mu_B B}{\pi\hbar c}
 =\frac{eB\lambda^2}{2\pi m_ec}
 \simeq35\,\mathrm{pm}}.
\]
\excheck 题目问最外侧两条之\emph{总间隔}，
不是中央至侧线的间隔。

\atomictermsproblem{161}{250}
光谱仪分辨本领$\lambda/\delta\lambda=10^5$，
所观察单重项跃迁谱线$\lambda=536\,\mathrm{nm}$，
观察方向垂直于磁场。
求刚能分辨塞曼线分量的最小$B$。
\exsolution
光谱仪可分辨的邻线间距
$\delta\lambda=\lambda/10^5$。
正常塞曼三线的\emph{相邻}线距
\[
 \Delta\lambda_{\rm adj}
 =\frac{\lambda^2}{2\pi c}
 \frac{\mu_BB}{\hbar}
 =\frac{eB\lambda^2}{4\pi m_ec}.
\]
令$\Delta\lambda_{\rm adj}=\delta\lambda$，
\[
 \boxed{B_{\min}=
 \frac{4\pi m_ec}{e\lambda(10^5)}
 \simeq0.40\,\mathrm T=4.0\,\mathrm{kG}}.
\]
\excheck 沿垂直磁场方向可见$\pi$及$\sigma$分量，
而沿磁场方向$\pi$分量不可观察；
本题采用前者的邻线分辨条件。

\atomictermsproblem{162}{250}
在弱磁场中，
\({}^{3}D_1\to{}^{3}P_0\)谱线发生塞曼劈裂。
垂直磁场观察，相邻谱线角频率间隔
$\Delta\omega=1.32\times10^{10}\,\mathrm{s^{-1}}$。
求磁感应强度$B$。
\exsolution
上态$L=2,S=1,J=1$，
朗德因子$g_u=1/2$；
下态$J_l=0$，不劈裂。
允许跃迁的上态$m_u=0,\pm1$，
相邻频差为$g_u\mu_BB/\hbar$，
因此
\[
 \boxed{B=\frac{\hbar\Delta\omega}{g_u\mu_B}
 =\frac{2\hbar(1.32\times10^{10})}{\mu_B}
 \simeq0.300\,\mathrm T=3.0\,\mathrm{kG}}.
\]
\excheck 要用上态$g_u=1/2$，
不是盲目套普通三重线$g=1$。

\atomictermsproblem{163}{250}
钠黄线双线波长分别为
$589.59$和$589.00\,\mathrm{nm}$。
(a)求弱场下\({}^{2}P_{3/2}\)与
\({}^{2}P_{1/2}\)两项相邻塞曼子级的能距比；
(b)当\({}^{2}P_{3/2}\)的相邻磁子能级间距
为自然精细分裂的$1/\eta$，$\eta=50$时，
求磁场$B$。
\exsolution
(a) 对$L=1,S=1/2$，
朗德因子分别
$g_{3/2}=4/3$、$g_{1/2}=2/3$，
相邻塞曼间距是$g\mu_BB$，
\[
 \boxed{\frac{\delta E_{3/2}}{\delta E_{1/2}}
 =\frac{4/3}{2/3}=2:1}.
\]
(b) 自然精细分裂
\[
 \Delta E_{\rm fs}=hc\left|
 \frac1{589.00\,\mathrm{nm}}-
 \frac1{589.59\,\mathrm{nm}}\right|.
\]
令
$g_{3/2}\mu_BB=\Delta E_{\rm fs}/50$，
得到
\[
 \boxed{B=\frac{\Delta E_{\rm fs}}
 {50(4/3)\mu_B}
 \simeq0.55\,\mathrm T=5.5\,\mathrm{kG}}.
\]
\excheck (a)的比为2而非钠两种波长的比，
因为决定邻级能距的是朗德因子。

\atomictermsproblem{164}{250}
绘出弱磁场中\({}^{2}P_{3/2}\to{}^{2}S_{1/2}\)
允许跃迁关系，并求$B=4.5\,\mathrm{kG}$下
全部塞曼分量的角频率偏移。
\exsolution
上态$J_u=3/2,g_u=4/3$，
$m_u=\pm3/2,\pm1/2$；
下态$J_l=1/2,g_l=2$，
$m_l=\pm1/2$。
只保留$\Delta m=m_u-m_l=0,\pm1$的六条：
\[
 \begin{array}{c|c|c}
 m_u&m_l&g_um_u-g_lm_l\\\hline
 +3/2&+1/2&+1\\
 +1/2&+1/2&-1/3\\
 +1/2&-1/2&+5/3\\
 -1/2&+1/2&-5/3\\
 -1/2&-1/2&+1/3\\
 -3/2&-1/2&-1
 \end{array}
\]
每条偏移为
$\Delta\omega=(g_um_u-g_lm_l)\mu_BB/\hbar$。
取$B=0.45\,\mathrm T$，
$\mu_BB/\hbar\simeq3.96\times10^{10}\,\mathrm{s^{-1}}$，
因此
\[
 \boxed{\Delta\omega=
 (\pm1.32,\ \pm3.96,\ \pm6.60)\times10^{10}\,
 \mathrm{rad/s}}.
\]
\excheck 六条分量包括$\Delta m=0$两条$\pi$线，
和$\Delta m=\pm1$四条$\sigma$线；
完整图的箭头关系由表逐行确定。

\atomictermsproblem{165}{250}
图6.9有磁极N、S及小反射镜M：
从方向1（与磁场垂直）及经镜面从方向2
（等效沿磁场方向）观察同一条反常塞曼谱线。
两方向分别能看到多少条分量：
(a)\({}^{2}P_{3/2}\to{}^{2}S_{1/2}\)；
(b)\({}^{3}P_2\to{}^{3}S_1\)？
\exsolution
塞曼电偶极选择规则$\Delta m=0,\pm1$。
垂直于外场方向观察时，
$\pi$（$\Delta m=0$）与$\sigma$（$\Delta m=\pm1$）
两类分量都能出现；
沿磁场方向观察时$\pi$分量不可见。
(a) 题164已经枚举全部六条允许跃迁，
其中两条$\Delta m=0$、四条$\Delta m=\pm1$，故
\[
 \boxed{\text{(a)方向1：6条；方向2：4条}}.
\]
(b) 上态$J=2$的$m=-2,-1,0,1,2$，
下态$J=1$的$m=-1,0,1$。
枚举$\Delta m=0$有3条，
$\Delta m=+1$有3条，
$\Delta m=-1$有3条，合计9条；
沿场方向只看见六条$\sigma$分量，故
\[
 \boxed{\text{(b)方向1：9条；方向2：6条}}.
\]
\excheck 源图6.9的镜面使方向2等效观察沿磁场传播的光；
不能简单因“镜面左右对称”而认为两方向分量数相同。

\atomictermsproblem{166}{250}
弱磁场$B=3.4\,\mathrm{kG}$下，
\({}^{3}D_3\to{}^{3}P_2\)
谱线的总塞曼角频率分裂宽度$\Delta\omega$为多少？
\exsolution
上态$S=1,L=2,J=3$，
$g_u=1+[12+2-6]/24=4/3$；
下态$S=1,L=1,J=2$，
$g_l=1+[6+2-2]/12=3/2$。
频移是
\[
 \delta\omega=(g_um_u-g_lm_l)\frac{\mu_BB}{\hbar},
 \quad |m_u-m_l|\le1.
\]
枚举允许磁量子数，
最大正频移为
$(5/3)\mu_BB/\hbar$，
最小为负对称值
$-(5/3)\mu_BB/\hbar$，
故端到端角频率全宽
\[
 \boxed{\Delta\omega=
 \frac{10}{3}\frac{\mu_BB}{\hbar}
 \simeq9.97\times10^{10}\,\mathrm{rad/s}
 \approx1.0\times10^{11}\,\mathrm{rad/s}}.
\]
\excheck “总分裂”指最外两条塞曼分量的角频率差，
不能只报最大单侧位移$5\mu_BB/(3\hbar)$。
% === END 36-atomic-spectra-159-166.tex ===
